\section{成败试验中的 Bernoulli、二项与等待分布}
\label{sec:05-Probability/06-Common-Distributions/01}

一个接口，每次调用要么返回成功要么超时。同一份调用日志摆在三个人面前，问出了三个问题。

运维想知道：今早那 \(10\) 次调用里，超时了几次？值班的人想知道：现在开始重试，第几次才能成功？做容量评估的人想知道：要凑够 \(3\) 次成功，总共得调用多少次？

三个问题共用一个超时率 \(p\)，共用同一条日志，答案却落在三个不同的分布里。分开它们的不是 \(p\)，是谁被你固定下来、谁留给随机性。

\begin{center}
\begin{tikzpicture}[x=1cm,y=1cm,font=\footnotesize,>=stealth]
  \foreach \i in {1,...,14}{\draw[black!45,fill=black!5] (0.62*\i-0.62,0) rectangle (0.62*\i,0.62);}
  \foreach \i in {3,7,8,12,13}{\draw[black!45,fill=blue!30] (0.62*\i-0.62,0) rectangle (0.62*\i,0.62);
    \node at (0.62*\i-0.31,0.31) {\(\checkmark\)};}
  \foreach \i in {1,...,14}{\node[black!55,font=\scriptsize] at (0.62*\i-0.31,-0.3) {\i};}
  \draw[black!60] (0,1.02) -- (0,1.17) -- (6.2,1.17) -- (6.2,1.02);
  \node[above,black!75] at (3.1,1.2) {固定 \(n=10\)：随机的是成功数 \(3\)};
  \draw[black!60] (0,-0.62) -- (0,-0.77) -- (1.86,-0.77) -- (1.86,-0.62);
  \node[right,black!75] at (2.0,-0.77) {等到第一次成功：随机的是试验数 \(3\)};
  \draw[black!60] (0,-1.32) -- (0,-1.47) -- (4.96,-1.47) -- (4.96,-1.32);
  \node[right,black!75] at (5.1,-1.47) {等到第 \(3\) 次成功：随机的是试验数 \(8\)};
\end{tikzpicture}

\emph{一条 0--1 序列，打勾的格子表示成功。三条括线是三种读法：上面那条圈住前十格，数里面有几个勾；下面两条从头出发，数到第一个勾、第三个勾各花了几格。序列没变，随机变量换了三次。}
\end{center}

上面这张图是本节的全部内容。剩下的工作是把三种读法各自的概率写出来，并且看清它们共用哪个零件。

\subsection{一次试验就是一个指示变量}

零件是 Bernoulli 分布。随机变量 \(X\) 只取 \(0\) 和 \(1\)，\(P(X=1)=p\)，\(P(X=0)=1-p\)。它简单到不像一个分布，但后面每个公式都由它拼出来。

因为 \(X\) 只取 \(0\) 和 \(1\)，有 \(X^2=X\)，于是 \(\E[X]=p\) 而 \(\Var(X)=p-p^2=p(1-p)\)。方差在 \(p=1/2\) 时最大，在两端趋于零：结果几乎确定的时候没什么可波动的。这条曲线在后面会反复出现——二项的方差、Beta 的方差、二元熵，形状都来自同一个 \(p(1-p)\)。

\subsection{固定次数，数成功：Binomial}

图中最上面那条括线。做 \(n\) 次独立试验，每次成功概率同为 \(p\)，记成功总数为 \(S_n\)。

为什么答案不是 \(p^n\)？因为 \(S_n=k\) 并不指定成功落在哪几次上。任何一条具体的成功--失败序列，只要含 \(k\) 个成功，概率都是 \(p^k(1-p)^{n-k}\)；而含 \(k\) 个成功的序列共有 \(\binom nk\) 条，彼此互斥。相加即得

\[
P(S_n=k)=\binom nk p^k(1-p)^{n-k},\qquad k=0,1,\ldots,n.
\]

组合数在这里数的是位置，指数项给的是任一指定模式的概率——这正是\hyperref[sec:04-Discrete-Mathematics/01-Combinatorics/02]{``二项式与双重计数''}里那个展开式的概率版本。

\(S_n\) 是 \(n\) 个独立 Bernoulli 之和，所以 \(\E[S_n]=np\)，\(\Var(S_n)=np(1-p)\)，不需要另做求和。

三条前提缺一不可，而且每条都对应一种现实中的失效。次数必须事先定死：中途因为``结果看起来不错''就收手，随机变量的含义就变了。每次的成功概率必须相同：把不同型号的零件、不同流量来源的请求混在一起统计，得到的是 Poisson--binomial 而不是 Binomial，方差会偏小。各次必须独立：从一批有限的对象里不放回地抽，前面的结果会改写后面的概率，这是下一节的主题。

\subsection{固定成功数，数试验：Geometric 与 Negative Binomial}

图中下面那两条括线。序列没变，只是这次不预先划定长度，而是等到某个事件出现为止。

等第一次成功：要让 \(T=k\)，前 \(k-1\) 次必须全失败、第 \(k\) 次成功，独立相乘即得 \(P(T=k)=(1-p)^{k-1}p\)。期望是 \(1/p\)，方差是 \((1-p)/p^2\)。期望的形状很好记：成功率 \(p=0.2\) 时平均五次里出一次，所以平均要等五次。这就是 Geometric 分布。

等第 \(r\) 次成功：要让 \(T_r=k\)，前 \(k-1\) 次里恰有 \(r-1\) 次成功、第 \(k\) 次成功，于是

\[
P(T_r=k)=\binom{k-1}{r-1}p^r(1-p)^{k-r},\qquad k=r,r+1,\ldots
\]

这是 Negative Binomial 分布。它的期望和方差不用重算：每次成功之后，等下一次成功的过程和从头开始等完全一样，所以 \(T_r\) 是 \(r\) 个独立 Geometric 之和，\(\E[T_r]=r/p\)，\(\Var(T_r)=r(1-p)/p^2\)。名字里的``负''来自 \((1-p)^{-r}\) 的级数展开系数，跟怎么理解它没什么关系。

这两个分布有一个共同的坑，值得单独提醒。有些文献让随机变量记录的不是试验总数，而是成功之前的失败次数，支撑从 \(0\) 起算，期望相应地变成 \((1-p)/p\) 与 \(r(1-p)/p\)。两套公式长得很像，代入的却是不同的量。拿到一段代码或一张表先确认变量的含义，比核对公式重要。

\subsection{等待有多久，与已经等了多久无关}

Geometric 分布满足一条别的离散分布都没有的性质：

\[
P(T>s+t\given T>s)=P(T>t).
\]

验证只要一步。生存函数是 \(P(T>k)=(1-p)^k\)（前 \(k\) 次全失败），代入条件概率的定义，\((1-p)^{s+t}/(1-p)^s=(1-p)^t\)。

读出来的意思是：已经连失败 \(s\) 次，接下来还要等多久，分布和一个刚开始等的人完全相同。历史在概率层面被抹掉了，没有任何``该轮到我了''的势头累积。

这是好消息也是坏消息。如果你模拟的确实是一串互不相干的独立尝试——理想的随机数、无状态的重试——它精确成立。而现实中的等待多半带记忆：设备会老化，队列会积压，重试会因为退避策略越隔越久，操作员会学习。用无记忆的模型去拟合有记忆的过程，中心附近可能看不出问题，尾部预测会明显偏。下一章的 Exponential 分布会把同一件事搬到连续时间上，那里的讨论更细。

\subsection{停止规则是模型的一部分}

回到最初那张图。四个分布的分工可以压成一句话：Bernoulli 看一次，Binomial 固定次数问成功数，Geometric 和 Negative Binomial 固定成功数问试验数。前者的 \(n\) 由你决定，后者的 \(n\) 由数据决定。

最常见的一类混淆是把``做了 \(n\) 次，成功 \(k\) 次''和``等到 \(k\) 次成功，做了 \(n\) 次''当成同一件事。两句话里都出现 \(n\)、\(k\) 和 \(p\)，但哪个被固定完全相反，随机变量也就完全不同。

这件事在实验分析里有直接后果。假设你在跑一个线上 A/B 测试，计划采集 \(10000\) 次曝光后再看结果，那么成功数服从 Binomial，围绕它做的检验是有效的。可如果你改成每天盯一次仪表盘，看到显著就停——你实际执行的是一条``等到差异够大为止''的停止规则，采集量本身成了随机变量。此时的样本量分布更接近等待型而非固定型，而按固定样本量算出来的 \(p\) 值会系统性偏小，把噪声读成效果。要么事先锁死样本量，要么改用为序贯观测设计的方法；把停止规则留在脑子里而不写进模型，是这类分析出错最常见的入口。

本节假定每次试验的成功概率始终是同一个 \(p\)，而且各次互不影响。下一节只拿走后面那一条：从一箱固定组成的对象里不放回地抽，前一次抽走了什么，会改写后一次的概率。
